\chapter{Style as an objective function}
\label{ch:style}

\begin{objectives}
\begin{itemize}
\item Build entropy from scratch and read it as "how many moves the opponent effectively has".
\item Compute policy entropy and perplexity on real positions.
\item Discover, from real data, why maximising opponent confusion is \emph{not} the same as
  creating danger --- and repair the objective.
\item Compare four ways to steer an engine, and rank them by how much they corrupt the
  evaluation.
\item State the constrained-selection pattern that gets style without giving up soundness.
\end{itemize}
\end{objectives}

\section{What "style" is, mathematically}

An engine plays the move that maximises something. Change what it maximises and you change
how it plays. That is the whole of engine steering, and stating it that plainly is useful,
because it forces the real question: \emph{what exactly do you want maximised?}

"Sharp, dynamic, sacrificial, Tal-like" is a feeling. To steer an engine you need a number.
This chapter builds candidate numbers, tests them against real output, and finds that the
obvious one is wrong in an instructive way.

\section{Entropy, from scratch}

\subsection{Surprise}

How surprised should you be by an event of probability $p$? Three requirements: a certain
event ($p=1$) is no surprise; rarer is more surprising; and the surprise of two independent
events together should be the \emph{sum} of their surprises. Exactly one function does all
three:
\begin{equation}
  \label{eq:surprise}
  \text{surprise}(p) \;=\; -\ln p .
\end{equation}
Check: $-\ln 1 = 0$ \checkmark; smaller $p$ gives bigger value \checkmark; and
$-\ln(p_1p_2) = -\ln p_1 - \ln p_2$ \checkmark.

\subsection{Entropy is average surprise}

Given a distribution $P$ over moves, its \textbf{entropy} is the surprise you should expect:
\begin{equation}
  \label{eq:entropy}
  \Ent(P) \;=\; \sum_a P(a)\cdot\bigl(-\ln P(a)\bigr) \;=\; -\sum_a P(a)\ln P(a).
\end{equation}

\begin{worked}{entropy of two policies}
\emph{(Illustrative.)} Position A: one move overwhelmingly indicated,
$P = (0.95, 0.03, 0.02)$:
\[
\Ent = -(0.95\ln0.95 + 0.03\ln0.03 + 0.02\ln0.02)
     = -(-0.0487 - 0.1052 - 0.0782) = 0.232 .
\]
Position B: three equally plausible moves, $P = (0.34, 0.33, 0.33)$:
\[
\Ent = -(0.34\ln0.34 + 0.33\ln0.33 + 0.33\ln0.33) = 1.098 \;\approx\; \ln 3 .
\]
Nearly five times the entropy of A.
\end{worked}

\subsection{Perplexity: entropy in units you can feel}

Entropy in nats is hard to have intuitions about. Exponentiate it:
\begin{equation}
  \label{eq:perplexity}
  \text{perplexity}(P) \;=\; e^{\Ent(P)}
\end{equation}
and you get \textbf{the effective number of options}. A distribution spread evenly over $k$
moves has perplexity exactly $k$; one concentrated on a single move has perplexity 1. Position
A above has perplexity $1.26$ ("basically one move"); position B has $3.00$ ("genuinely three").

\begin{realdata}[title={Real engine output: policy perplexity of this book's positions}]
\begin{center}
\begin{tabular}{@{}lrrrl@{}}
\toprule
Position & legal & $\Ent(P)$ & perplexity & reading \\
\midrule
K+P, Black to move & 2 & $0.689$ & $1.99$ & both moves live \\
K+P, White to move & 4 & $1.032$ & $2.81$ & effectively three \\
Drawn opposition ending & 5 & $1.609$ & $\mathbf{5.00}$ & \textbf{all five equally plausible} \\
Initial position & 20 & $2.420$ & $11.25$ & eleven real options \\
Morphy, before O-O-O & 49 & $2.506$ & $12.26$ & \\
Morphy, before Qb8+ & 46 & $2.618$ & $13.71$ & \\
\bottomrule
\end{tabular}
\end{center}
\end{realdata}

\section{The obvious objective, and why it is wrong}

Tal's chess is usually described as making the opponent's task hard: many plausible moves,
most of them losing. The natural formalisation --- and the one in this project's earlier
theory notes --- is to prefer moves that maximise the \emph{opponent's} policy entropy:
\begin{equation}
  \label{eq:talnaive}
  \text{Score}_{\text{Tal}}(s') \;=\; V(s') \;+\; \gamma\,\Ent\bigl(P(\cdot\given s')\bigr) .
\end{equation}
Evaluate each candidate move, look at the position it produces, and reward those in which the
opponent's policy is flat.

Now look at the third row of the table above.

\begin{pitfall}
\textbf{The drawn king-and-pawn ending has the maximum possible policy entropy.} Perplexity
$5.00$ out of 5 legal moves --- the opponent's policy is perfectly flat, more confused than in
any other position in our data set, including a position with mate in two on the board.

And the position is a dead draw. Stockfish at depth 30 gives $0.00$ for every move; Leela's
draw probability is $0.997$. The opponent has five equally plausible moves \emph{because
nothing matters}. Every move draws, so the policy has no reason to prefer any of them.

Equation~\eqref{eq:talnaive} would steer you here. It scores maximum "Tal" on the most
sterile position in the book.
\end{pitfall}

This is worth dwelling on because it is a real defect in a real design note, found by
computing the metric on real positions rather than reasoning about it. Entropy measures
\emph{indifference}, and indifference has two utterly different causes:

\begin{center}
\begin{tabular}{@{}lll@{}}
\toprule
 & many plausible moves because\ldots & \\
\midrule
\textbf{Minefield} & several moves are critical and only one holds & what you want \\
\textbf{Swamp} & nothing anyone does matters & what you get \\
\bottomrule
\end{tabular}
\end{center}

\subsection{The repair}

The missing ingredient is \emph{stakes}, and Chapter~\ref{ch:currency} already built it:
sharpness $= 1 - d$, the probability that the game gets decided. Multiply rather than add,
so that zero stakes zeroes the whole term:
\begin{equation}
  \label{eq:talfixed}
  \boxed{\;
  \text{Score}(s') \;=\; V(s') \;+\; \gamma\;\underbrace{\bigl(1 - d(s')\bigr)}_{\text{is anything at stake?}}
  \cdot\underbrace{\Ent\bigl(P(\cdot\given s')\bigr)}_{\text{how many ways to go wrong?}} \;}
\end{equation}

Recompute the drawn ending: $1 - d = 0.003$, so the styling term is
$0.003 \times 1.609 = 0.005$ regardless of $\gamma$. The swamp is correctly ignored. A
position with $d = 0.15$ and the same entropy scores $0.85 \times 1.609 = 1.37$ --- nearly
three hundred times more.

\begin{keyidea}
\textbf{A steering term needs both a "how confusing" factor and a "does it matter" factor,
multiplied.} Either alone is trivially gameable: entropy alone steers into dead positions;
stakes alone steers into any decisive position including simple ones. This is the general
shape of every honest complexity metric.
\end{keyidea}

\begin{worked}{choosing a steer move from real data}
The K+P ending, using the opponent's policy after each of White's four moves (measured):
\begin{center}
\begin{tabular}{@{}lrrrrl@{}}
\toprule
White's move & opp.\ legal & $\Ent(P_{\text{opp}})$ & perplexity & $V$ & Stockfish \\
\midrule
Kd6 & 3 & $1.066$ & $2.90$ & $+0.968$ & mate in 11 \\
Kf6 & 3 & $1.023$ & $2.78$ & $+0.982$ & mate in 19 \\
Kd5 & 5 & $0.967$ & $2.63$ & $0.000$ & draw \\
Kf5 & 5 & $0.959$ & $2.61$ & $0.000$ & draw \\
\bottomrule
\end{tabular}
\end{center}
Apply Equation~\eqref{eq:talfixed} with $\gamma = 0.1$: the two drawing moves are excluded by
$V$; between the two winning moves, Kd6 leaves Black with $2.90$ effective options against
Kf6's $2.78$, so Kd6 is both objectively better \emph{and} the one that gives the defender
more chances to go wrong. Here style and truth agree, which is the easy case.

The interesting case is when they do not, and that is exactly what §16.5 is about.
\end{worked}

\section{Four levers, ranked by how much damage they do}

\subsection{Lever 1: bias the search's $Q$}

\[
  Q_{\text{biased}}(s,a) \;=\; Q(s,a) \;+\; \lambda\,M(\text{motif along the line}) .
\]
Add a motif bonus inside the search. This is the most direct and the most dangerous: the
biased value propagates up the tree, so the engine's \emph{evaluation} is now a blend of
"winning" and "stylish", and you can no longer tell which. Worse, it invites the horizon
effect --- the engine will play a real material loss to reach a position that scores the
bonus, because the bonus is inside the quantity being maximised.

\begin{pitfall}
This lever also destroys the thing that makes your tool trustworthy. Once $Q$ is
style-adjusted, every displayed evaluation, every "you lost $0.4$ here", and every drill
threshold is measured in a corrupted currency. \textbf{Never bias the number you also report.}
\end{pitfall}

\subsection{Lever 2: an auxiliary head with a dial}

Train the network to emit a second value $V_{\text{style}}$ alongside $V_{\text{win}}$, and
blend at inference:
\[
  V_{\text{steered}} = (1-\alpha)V_{\text{win}} + \alpha V_{\text{style}} .
\]
Clean in principle and it keeps the two quantities separable --- but it requires training a
network, which is out of scope here by several orders of magnitude of compute. Worth knowing
the option exists; not worth planning around.

\subsection{Lever 3: re-rank the candidates after search (the safe one)}

Do not touch the search at all. Let it run honestly, then choose among its candidates by a
style criterion, subject to a soundness constraint:
\begin{equation}
  \label{eq:constrained}
  \boxed{\;
  a_{\text{steer}} \;=\; \argmax_{a\,:\,Q(s,a)\;\ge\; Q^{*} - \epsilon}\;\; \text{Style}(s\!\cdot\!a)
  \;}
\end{equation}
where $Q^{*}$ is the best value found and $\epsilon$ is how much objective quality you are
willing to spend --- a number you choose, state, and can show the user.

\begin{keyidea}
Equation~\eqref{eq:constrained} is the pattern to build on. It has four properties the others
lack: the evaluation stays honest; the cost of style is \emph{explicit} ($\epsilon$, in units
the user understands); it cannot produce an unsound move, because unsound moves fail the
constraint; and it works with the engine you already have, today, with no training.
\end{keyidea}

\subsection{Lever 4: restrict the opponent (the anaconda)}

For the opposite style --- squeezing rather than storming --- the natural quantity is the
opponent's mobility:
\[
  R_{\text{mobility}}(s') \;=\; -\lambda\,\bigl|\mathcal{M}_{\text{opp}}(s')\bigr| ,
\]
the number of legal moves available to them. Cheap to compute (a move generator call) and it
plugs into Equation~\eqref{eq:constrained} in place of Style with no other changes.

It has the same failure mode as entropy, in mirror image: restricting a helpless opponent in a
dead position scores well and means nothing. It needs the same $(1-d)$ factor. In fact:

\begin{keyidea}
Every style metric in this chapter is of the form
\[
  \text{Style}(s') \;=\; \bigl(1 - d(s')\bigr) \times \bigl(\text{something about the
  opponent's options}\bigr) .
\]
The left factor is what makes it about chess; the right factor is what makes it about style.
Get the left factor wrong and you have built a metric that scores dead positions highly, which
is precisely the failure this project has hit before in a different guise.
\end{keyidea}

\section{Notes on evaluating a steer}

Three cheap disciplines, before this becomes a feature:

\begin{enumerate}
\item \textbf{Compute it on the swamp.} Run any candidate metric on the drawn opposition
  ending. If it scores well there, the metric is broken. Keep a small set of such
  \emph{negative witnesses} and run every metric against them --- this chapter found a real
  defect with a single position.
\item \textbf{Report $\epsilon$ in the interface.} "This move gives up $0.08$ of evaluation to
  reach a much sharper position" is honest, checkable and teaches the user something. "Steer
  move" alone is not.
\item \textbf{Measure whether the opponent actually goes wrong.} The whole justification is
  practical: that a human at the other end errs more often. That is testable against your own
  corpus --- do positions with high $\text{Style}$ actually precede more errors? If not, the
  metric is measuring something real but useless.
\end{enumerate}

\begin{repobox}[title={in this repository}]
The existing TS2 steer score already has components along these lines (decisiveness,
narrowness, policy trap, attention). The two contributions of this chapter are: the
multiplicative $(1-d)$ gate, which the entropy component needs and which the drawn-ending
witness shows is not optional; and the constrained form of Equation~\eqref{eq:constrained},
which makes the price of steering explicit rather than baked into a weighted sum.
\end{repobox}

\begin{exercises}

\exhead[surprise]{ch16:surprise} Compute $-\ln p$ for $p = 0.9, 0.5, 0.1, 0.01$. Why must
surprise be additive for independent events, and which other candidate function does that rule
out?

\exhead[entropy by hand]{ch16:entropy} Compute $\Ent$ and perplexity for $P = (0.7,0.2,0.1)$
and for $P = (0.4,0.35,0.25)$. Which position offers the opponent more real choice, and by how
many "effective moves"?

\exhead[maximum entropy]{ch16:maxent} Prove (or argue convincingly) that entropy over $k$
options is maximised by the uniform distribution, with value $\ln k$. Verify against the
drawn ending's measured $1.609$ and $\ln 5$.

\exhead[the swamp]{ch16:swamp} In your own words, explain why the drawn ending has maximum
policy entropy. Then explain why this is a fact about the \emph{policy head's training
objective} rather than a quirk.

\exhead[repairing the objective]{ch16:repair} Compute the style term of
Equation~\eqref{eq:talfixed} for all six positions in the §16.2.3 table, using their real $d$
values from Chapter~\ref{ch:currency}. Rank them. Does the ranking match your chess intuition
about which is the most dangerous position to defend?

\exhead[additive versus multiplicative]{ch16:multiplicative} Why does
Equation~\eqref{eq:talfixed} multiply the two factors instead of adding them? Construct a
position pair where the additive version gives the wrong ranking.

\exhead[horizon effect]{ch16:horizon} Describe concretely how lever 1 could make an engine
sacrifice a rook for nothing. Which term in $Q_{\text{biased}}$ is responsible, and why does
increasing search depth not fix it?

\exhead[choosing epsilon]{ch16:epsilon} You must choose $\epsilon$ in
Equation~\eqref{eq:constrained}. Argue for a specific value for (a) a training tool for a
2100-rated player, (b) an engine playing a rated game. What information would let you choose
better than by taste?

\exhead[mobility]{ch16:mobility} Compute the opponent's mobility for each of the four moves in
the K+P ending from the data in §16.3.1 (the "opp.\ legal" column). Which move minimises it?
Is that the move you want, and what does the answer say about mobility as a standalone
metric?

\exhead[a negative witness set]{ch16:witness} Propose five positions --- described in words ---
that any style metric must score \emph{low}. For each, say which naive metric it would catch
out.

\exhead[the real test]{ch16:realtest} Design the experiment that would tell you whether your
steer metric predicts human error, using your 9{,}030-game corpus. What would you compute,
what would the null result look like, and what confound would you have to control for?

\end{exercises}
